\section{The agent-facing registry API}\label{sec:agent}

\statspai{} exposes the same surface used by a human Python user to
LLM agents through typed schemas and an MCP server. The paper's claim is
mechanical and contractual: schemas, typed errors, result handles,
citations, and deterministic traces make machine-orchestrated calls
inspectable by the same validation ledger as human calls. It is not a
behavioural claim that an autonomous agent will always choose the right
empirical design. The agent layer is a software-interface contribution
rather than an intelligence benchmark: it reduces four otherwise opaque
handoffs -- tool discovery, failure recovery, result reuse, and citation
resolution -- to artifacts a reviewer can inspect in the same archive,
which is why the validation target is schema/MCP/citation
reproducibility, not a model-quality score.

\subsection{Why a schema, not a docstring}\label{sec:why-schema}

Docstrings, \proglang{R} help pages, and \proglang{Stata} help files are
prose interfaces. A tool-using agent needs a typed object containing
parameter names, types, defaults, required fields, enumerated choices,
and a short description including stability and limitations when
relevant. The schema call returns the executable OpenAI-style object;
the description and agent-card calls return richer planning metadata
such as assumptions, preconditions, failure modes, alternatives, typical
sample size, references, and examples.

\subsection{The function-schema layer}\label{sec:fn-schema}

The discovery surface lists registered functions, describes a selected
function, emits its executable schema, and returns curated agent cards.
These functions are themselves registered, which lets an agent inspect
the available estimator catalogue before selecting a method.
Complementary design-detection and preflight calls expose cheap planning
diagnostics before an estimator is run.

\subsection{The typed exception envelope}\label{sec:exception-envelope}

Estimator failures inherit from \code{StatsPAIError} and carry
structured fields: \code{error\_kind}, \code{detail},
\code{recovery\_hint}, \code{alternative\_functions}, and any partial
\code{diagnostics}. Warnings use the same idea for non-fatal
degradations. The MCP server serializes these records as structured
errors, so an agent can branch on \code{weak\_instrument},
\code{no\_overlap}, \code{rank\_deficient}, or similar identifiers
without scraping English messages.

\subsection{The estimand-first DSL and robustness audit}\label{sec:causal-question}

\code{sp.causal\_question(...)} accepts a requested estimand, treatment,
outcome, data, optional design, covariates, instruments, and sensitivity
flag, then routes to the appropriate estimator and returns a
\code{CausalQuestionResult}. The result bundles the fitted estimator,
the design-detection provenance, optional sensitivity output, and an
\code{audit\_report}. The standalone \code{sp.audit(result)} call
produces the same missing-evidence checklist for any fitted result and
points to follow-up functions such as RD density tests, Honest-DiD
sensitivity, or unobserved-confounding bounds. Each check carries one of
four statuses --- passed, failed, missing, or not applicable --- and the
last is decided from the fit's own metadata by an explicit
applicability rule (an over-identification test is not applicable to a
just-identified IV model, and says why), so an agent is never sent after
a test that does not exist. The companion \code{sp.validation\_scope}
answers the neighbouring question of which validation artifacts cover
the fitted configuration (Section~\ref{sec:validation-tiers}).

\subsection{The bibliography verifier}\label{sec:cite}

Every public-facing citation emitted by the package must come from the
curated \code{paper.bib}. Function specs store bibliography keys, result
objects expose \code{cite(format=...)}, and tests assert that registered
keys parse cleanly and do not collide on author/year/title triples.
When a key is missing, the package raises an error rather than emitting
a plausible-looking citation. This targets a concrete agent failure
mode: hallucinated references.

\subsection{The MCP server}\label{sec:mcp}

The \code{statspai-mcp} server exposes registered functions as MCP
tools. Tool inputs are generated from the registry schemas, outputs are
JSON-safe result payloads plus \code{result\_id} handles, and typed
errors use the exception envelope above. Chained tools cover auditing,
Honest-DiD sensitivity, generic sensitivity analysis, CS aggregation,
Bacon decomposition, paper generation, and citation, and they accept
result handles instead of requiring an agent to marshal coefficient,
covariance, or cross-fit arrays through model context. The
deterministic trace in Section~\ref{sec:agentbench} exercises this
contract.

\subsection{Comparison with existing platforms and prior art}\label{sec:agent-compare}

Generating a JSON schema from a typed signature is solved by
\pkg{Pydantic}, \pkg{FastAPI}, and the LLM-vendor tool adapters; what
is package-specific here is that the
schema is the statistical registry itself, tied to validation status
and typed statistical errors, with result handles for
estimate-to-audit-to-sensitivity chains -- a validation-tiered tool
catalogue with a citation verifier, which statistical help systems do
not expose.
